\chapter*{Sommario}
\addcontentsline{toc}{chapter}{Sommario}
\markboth{Sommario}{Sommario}

Il \emph{quantum machine learning} ibrido propone di sostituire strati parametrici
classici con circuiti quantistici variazionali (VQC) all'interno di reti addestrate
per retropropagazione. La letteratura attribuisce a questa sostituzione quattro
benefici distinti: uno speedup computazionale, una
maggiore espressività a parità di parametri, una regolarizzazione implicita
indotta da unitarietà e limitatezza dell'osservabile misurato, e una migliore
separabilità delle classi in spazi di feature ad alta dimensione.

Questa tesi valuta sperimentalmente i tre benefici che non richiedono hardware
quantistico, lasciando esplicitamente fuori portata lo speedup. Tutti gli
esperimenti girano su un simulatore classico di vettore di stato, il cui costo
cresce esponenzialmente con il numero di qubit. Le misure riguardano circuiti
piccoli simulati e non confrontano complessità o tempi con hardware quantistico.
Non possono quindi dimostrare alcun vantaggio computazionale.

Sono formulate quattro ipotesi operative, ciascuna con metrica primaria,
direzione attesa e criterio di decisione dichiarati a priori, e valutate
con quattro esperimenti indipendenti: proiezioni $Q,K,V$ quantistiche in un
Transformer autoregressivo; \emph{patch embedding} quantistico in un Vision
Transformer; uno studio di ablazione a tre varianti sulla regolarizzazione
implicita; un confronto fra kernel di fedeltà quantistica e similarità coseno.
Il metodo sperimentale usato impiega baseline classiche appaiate, contabilizza gli squilibri
di parametri e, nell'ablazione sulla regolarizzazione, usa un controllo classico
con più parametri del blocco quantistico; adotta inoltre un design multi-seed e
multi-dataset e un'analisi statistica basata su differenze appaiate, intervalli
di confidenza parametrici e bootstrap, test di Wilcoxon e dimensione dell'effetto.

Nelle condizioni testate nessuna delle quattro ipotesi è confermata nella
direzione attesa. La perplexity aumenta di $0{,}067$ (IC 95\%
$[-0{,}141;0{,}275]$); nel Vision Transformer l'accuracy diminuisce di $0{,}024$
(IC $[-0{,}054;0{,}006]$) e la macro-AUC di $0{,}010$ con IC interamente
negativo. Nell'ablazione la differenza di gap fra variante quantistica e baseline
classica a output limitato è positiva su tutti i dataset ($+0{,}148$, $+0{,}100$,
$+0{,}085$), mentre il kernel IQP ha
un criterio di Fisher inferiore sia al coseno sia all'RBF.
Il costo di simulazione misurato è di uno-due ordini di grandezza superiore a
quello delle baseline classiche.

Oltre ai risultati sperimentali, il lavoro produce
\code{quantum\_framework}, una libreria per rendere ripetibile il confronto
quantistico/classico appaiato. Comprende layer variazionali, caricamento dei
dati, ciclo di addestramento con diagnostica su gradienti e saturazione,
metriche, test statistici e manifest di provenienza. Il protocollo sperimentale
documentato consente a terzi di replicare o confutare le affermazioni riportate.

Questo elaborato descrive un lavoro svolto in collaborazione con il collega
\Coautore{}.

\cleardoublepage
